\documentclass[a4paper,11pt]{article}

\usepackage[ngerman]{babel}
\usepackage[top=3cm,bottom=3cm,left=3cm,right=3cm]{geometry} %NICHT IN KEMIE2


%\usepackage{microtype} %schöneres typesetting  mit [stretch=10] für nur 1% anstelle von 2%
%\usepackage{lmodern} %Latin Modern FONT
%\usepackage{pifont} %schönere Font



\usepackage{amsmath}
\usepackage{nicefrac}
\usepackage{amssymb}
%\usepackage{mathtools}

\usepackage{titlesec} %Um Section, etc bearbeiten zu können
\usepackage{setspace}
\usepackage{enumitem} %enumerate
\usepackage{multicol}
\usepackage[many]{tcolorbox}
\usepackage{hyperref}
\usepackage{arydshln}

%\renewcommand{\ttdefault}{lmtt} %Schriftart wird geändert zu Latin Modern Typewriter


\hypersetup{hidelinks,
			colorlinks=true,
			linkcolor=black,
			citecolor=miudiutex,
			urlcolor=miugray2
			}

\setlength{\parindent}{0pt}
\setcounter{section}{0}


\titlespacing*{\section}{0pt}{20mm}{6mm}
\titlespacing*{\subsection}{0pt}{6mm}{3mm}
\titlespacing*{\subsubsection}{0pt}{3mm}{0mm}



%\definecolor{miudiutex}{HTML}{2f3d39}
\definecolor{miudiutex}{HTML}{1d2926}
\definecolor{miugray}{HTML}{b4cccf}
\definecolor{miugray2}{HTML}{4d7365}
\definecolor{red}{HTML}{cf1f5c}
\definecolor{green}{HTML}{00A36C}
\definecolor{delphinidin}{HTML}{315794}
\definecolor{pelargonidin}{HTML}{b33807}
\definecolor{cyanidin}{HTML}{ba0746}
\definecolor{petunidin}{HTML}{b56fed}



\raggedcolumns %wichtig für Multicolumns


%================================
% SYMBOLS
%================================

\newcommand{\RA}{$\rightarrow$~}
\newcommand{\RL}{$\rightleftharpoons$~}
\newcommand{\HARP}{\ding{226}~}
%$\xrightarrow{2}$ %Pfeil mit 2 drüber


%================================
% SHORT COMMANDS (BOXES ; ETC)
%================================


\newcommand{\notiz}[1]{\begin{notizz}{}{}\begin{center}
			{\color{miugray2!50!miudiutex}#1} \end{center}\end{notizz}}
\newcommand{\zit}[2]{\begin{zitat*}{#1}{}\small #2 \end{zitat*}}

\newcommand{\frage}[2]{\begin{qst*}{#1}{}\small #2 \end{qst*}}

\newcommand{\unten}{\end{center} \tcblower \begin{center}}


\newcommand*\circled[1]{\tikz[baseline=(char.base)]{
            \node[shape=circle,draw,inner sep=0.5pt,miudiutex] (char) {\tiny #1};}}
            
\newcommand{\miusquare}{{\color{miugray2!50}\scriptsize $\blacksquare$}~}
\newcommand{\miusquarp}{{\color{petunidin!50}\scriptsize $\blacktriangleright$}~}
\newcommand{\niu}{\item[\miusquare]}
\newcommand{\nip}{\item[\miusquarp]}

%%%%% ABSAETZE

\newcommand{\pps}{\par \vspace*{1mm}}
\newcommand{\pp}{\par \vspace*{3mm}}
\newcommand{\ppbig}{\par \vspace*{8mm}}
\newcommand{\pphuge}{\par \vspace*{10mm}}

%================================
% FRAGE BOX
%================================

\tcbuselibrary{theorems,skins,hooks}
\newtcbtheorem{qst}{Frage:}
{%
	enhanced,
	breakable,
	colback = gray!15!white,
	frame hidden,
	boxrule = 0sp,
	borderline west = {2.5pt}{0pt}{red!70},
	sharp corners,
	detach title,
	before upper = \tcbtitle\par\smallskip,
	coltitle = black,
	fonttitle = \bfseries\sffamily,
	description font = \mdseries,
	separator sign none,
	segmentation style={solid, gray!50!black},
}
{th}




%================================
% ZITAT BOX
%================================

\tcbuselibrary{theorems,skins,hooks}
\newtcbtheorem{zitat}{Zitat}
{%
	enhanced,
	breakable,
	colback = gray!15!white,
	frame hidden,
	boxrule = 0sp,
	borderline west = {2.5pt}{0pt}{black!70},
	sharp corners,
	detach title,
%	before upper = \tcbtitle\par\smallskip,
	coltitle = gray!50!black,
	fonttitle = \bfseries\sffamily,
	description font = \mdseries,
%	separator sign none,
	segmentation style={solid, gray!50!black},
}
{th}



%================================
% SIMPLE SCHWARZER RAHMEN BOX
%================================


\newcommand{\boxfr}[1]{
	\begin{tcolorbox}[
	drop shadow southeast,
	enhanced,
	colframe=black!50,
	colback=white,
	boxrule = 1pt, 
%	toprule=0pt, bottomrule=0pt,rightrule=0pt,leftrule=0pt,
	boxsep=5pt,
	arc=1pt,
	]
	\begin{center}
	#1
	\end{center}
	\end{tcolorbox}
}

\definecolor{miudiutex}{HTML}{1d2926}
\definecolor{miugray}{HTML}{b4cccf}
\definecolor{miugray2}{HTML}{4d7365}
\definecolor{white2}{HTML}{f5f7f8}
\definecolor{red}{HTML}{cf1f5c}
\definecolor{green}{HTML}{00A36C}
\definecolor{delphinidin}{HTML}{315794}
\definecolor{uniblau}{HTML}{004e9f}
\definecolor{unigelb}{HTML}{fcba00}
\definecolor{unigrau}{HTML}{909085}
\definecolor{pelargonidin}{HTML}{b33807}
\definecolor{cyanidin}{HTML}{ba0746}
\definecolor{paeonidin}{HTML}{d15ca6}
\definecolor{petunidin}{HTML}{b56fed}
\definecolor{malvidin}{HTML}{8a2d8c}

\newcommand{\metermilli}{\mskip3mu \text{mm}}
\newcommand{\cm}{\mskip3mu \text{cm}}
\newcommand{\m}{ \text{m}}
\newcommand{\lite}{ \text{l}}
\newcommand{\kg}{ \text{kg}}
\newcommand{\g}{ \text{g}}
\newcommand{\km}{ \text{km}}
\newcommand{\h}{ \text{h}}
\newcommand{\s}{ \text{s}}
\newcommand{\tonne}{ \text{t}}
\usepackage[utf8]{inputenc}
\usepackage[T1]{fontenc}
\usepackage{lmodern}

\usepackage[
    activate={true,nocompatibility},
    final,
    tracking=true,
    kerning=true,
    spacing=true,
    factor=1100,
    stretch=30,
    shrink=20
]{microtype}
\usepackage{pdfpages}
\usepackage{awesomebox}
\usepackage{marginnote}
\tcbuselibrary{documentation}
\usepackage{sidenotes}
\usepackage{import}
\usepackage{xifthen}
\usepackage{transparent}
\usepackage[table]{xcolor}


\newcommand{\abbicon}{%
  \protect\tikz[baseline=0.2mm,scale=0.8]{%
    % Das blaue Quadrat mit abgerundeten Ecken
    \fill[uniblau, rounded corners=1.2pt] (0,0) rectangle (0.8em, 0.8em);
    % Der weiße Kreis in der Mitte
    \fill[white] (0.4em, 0.4em) circle (0.12em);
  }%
} 
\usepackage[labelfont={bf},font={color=miudiutex,small},figurename={\abbicon $~$Abbildung}]{caption}


\newcommand{\bckgrnd}{
\AddToHookNext{shipout/background}{
\begin{tikzpicture}[remember picture, overlay]
 \draw[draw=none,fill=uniblau!70, shift={(current page.north west)}] (-3,0) rectangle (23,-3.75);
 \draw[draw=none,fill=miugray, shift={(current page.north east)}] (-7,0) rectangle (3,-3.75);
\end{tikzpicture}
}
}



\newcommand{\incfig}[2]{%
\def\svgwidth{#2\columnwidth}
\import{C://LaTeX-Files/pngs/}{#1.pdf_tex}
}

\renewcommand{\textbf}[1]{\textsf{{\bfseries {#1}}}}



\titleformat{\section}[hang]
		{\LARGE \color{uniblau!75!black} \sffamily \bfseries}
		{\thesection}
		{6mm}
		{}
		[]
\titleformat{\subsection}[hang]
		{\large \sffamily \bfseries \color{black}}
		{\thesubsection}
		{4mm}
		{{\color{miugray}$\blacksquare ~$}}

\titleformat{\subsubsection}[block]
		{\normalsize \bfseries \sffamily \color{miudiutex}}%
		{\normalsize \color{black} \thesubsubsection}
		{2mm}
		{{\color{miugray}$\blacksquare ~$}}
		[ \color{black}]
		
\titleformat{\paragraph}
		{\normalsize \bfseries \sffamily \color{black}}
		{\footnotesize \theparagraph}
		{1mm}
		{}

\pagestyle{empty}
\titlespacing*{\section}{0pt}{20mm}{12mm}
\titlespacing{\section}{0pt}{20mm}{12mm}

\titlespacing*{\subsection}{0pt}{14mm}{4mm}
\titlespacing{\subsection}{0pt}{14mm}{4mm}

\titlespacing*{\subsubsection}{0pt}{6mm}{3mm}
\titlespacing{\subsubsection}{0pt}{6mm}{3mm}

\titlespacing*{\paragraph}{0pt}{6mm}{2mm}
\titlespacing{\paragraph}{0pt}{6mm}{2mm}


\let\oldmarginfigure\marginfigure
\let\endoldmarginfigure\endmarginfigure

\let\oldmarginfigure\marginfigure
\let\endoldmarginfigure\endmarginfigure

\renewenvironment{marginfigure}[1][0pt] % [1] steht für 1 Argument, [0pt] ist der Standardwert
  {\oldmarginfigure[#1]\captionsetup{labelfont={sf,bf,color=uniblau!75!black},font={color=miudiutex,small}}}
  {\endoldmarginfigure}
  
  
  
  
\renewcommand{\labelitemi}{\color{uniblau!70}$\bullet$}
\renewcommand{\labelitemii}{\color{uniblau!70}$\cdot$}
\newcommand{\plmtsquare}{{\color{uniblau!70}$\blacksquare~$}}
\newcommand{\splmt}[1]{{\color{uniblau!75!black} \sffamily #1}}
\newcommand{\fplmt}[1]{{\color{uniblau!75!black} \sffamily \bfseries #1}}

\newcommand{\antwort}[1]{
\begin{tcolorbox}[
	enhanced,
	enforce breakable,
	standard jigsaw,
	boxrule=0.7pt,
	colframe=black,
	sharp corners,
	boxsep=5pt,
	title=\bfseries \sffamily Antwort,
	opacityback=0,
	coltitle=miudiutex,
	attach boxed title to top text left={yshift=-\tcboxedtitleheight/2},
	boxed title style={colback=white,colframe=white},
	bottomrule at break=0pt,
	toprule at break=0pt,
	pad at break=16mm,
]
	\begin{center}
	#1
	\end{center}
	\end{tcolorbox}
}




\rowcolors{2}{miugray!50}{miugray!50}
\arrayrulecolor{white} % Weiße Gitterlinien
\setlength{\arrayrulewidth}{1.5pt} % Etwas dickere Linien
\renewcommand{\arraystretch}{1.4} % Mehr Zeilenhöhe für bessere Lesbarkeit}


%\setlist[itemize]{itemsep=0mm}
\setlist[enumerate]{itemsep=0mm}




\newif\ifshowme
\showmetrue




\begin{document}
\newgeometry{top=1.8cm,bottom=4cm,left=1.5cm,right=7.5cm,marginparsep=-2.00cm,marginparwidth=4.75cm}


\bckgrnd


\begin{center}
\LARGE
\color{white}
\textbf{Übung 11} \par \vspace*{2mm}
\large \color{miugray} \textbf{Prozessbezogene Lebensmitteltechnologie}
\end{center}
\par 
\vspace*{6mm}


\color{miudiutex}
\section*{Mischungen \& Diffusion}

\subsection*{Aufgabe 1}
\marginnote[]{\small \fplmt{Anmerkung:} Ein Grad Brix entspricht der Dichte einer Saccharoselösung mit 1 g Saccharose
pro 100 g Lösung}
Aus einem Orangensaftkonzentrat mit 30 °Brix soll ein Orangensaft mit 11,5 °Brix hergestellt
werden. Wieviel kg Saft kann hergestellt werden, wenn Sie 5 kg Orangensaftkonzentrat zur
Verfügung haben?



\ifshowme
\antwort{
$m_1 = 5~\text{kg};~m_2 = m_1 + m_3 ;~m_3 = ~?~;~x_1=30~\text{°};~x_2=11,5~\text{°};~x_3=0~\text{°}$
\begin{align*}
m_1 \cdot x_1 + m_3 \cdot x_3 &= m_2 \cdot x_2 \\
m_1 \cdot x_1 + m_3 \cdot x_3 &= m_1 \cdot x_2 + m_3 \cdot x_2 \\
m_1 \cdot x_1 + m_3 \cdot x_3 - m_1 \cdot x_2 &= m_3 \cdot x_2 \\
m_1 \cdot x_1 - m_1 \cdot x_2 &= m_3 \cdot x_2 - m_3 \cdot x_3\\
m_1 \cdot x_1 - m_1 \cdot x_2 &= m_3 \cdot (x_2 - x_3)\\
\frac{m_1 \cdot x_1 - m_1 \cdot x_2}{x_2 - x_3} &= m_3 = 8,04
\end{align*}
Wir können also zu unseren vorhandenen 5 kg Orangensaftkonzentrat weitere 8,04 kg Wasser hinzugeben um einen °Brix von 11,5 zu erreichen. Dementsprechend erhalten wir 13,04 kg Orangensaft.

}
\pagebreak
\else
\fi


\subsection*{Aufgabe 2}
\marginnote[]{\small \fplmt{Anmerkung:} STP steht für Standard Temperature and Pressure, die Angabe bezieht sich also auf das Volumen bei Standardbedingungen.}
Ein Lebensmittel ist in Stickstoff-Schutzgasatmosphäre verpackt. Aufgrund eines Fehlers in
der Siegelnaht kommt es zur Diffusion von Sauerstoff und Wasserdampf aus der
Umgebungsluft in die Verpackung durch eine Pore (Porenradius: 12,5 µm, Porenlänge: 40
µm). Berechnen Sie den Stoffstrom für Sauerstoff und den Stoffstrom für Wasserdampf unter
der Annahme, dass die Konzentrationsdifferenzen von Sauerstoff und Wasserdampf 0,191 cm$^3$ (STP)/cm$^3$ und 12,91 µg/cm$^3$ betragen. Nehmen Sie für beide Gase einen Diffusionskoeffizienten von 0,243 cm$^2$/s an. 




\ifshowme
\antwort{
\textbf{Sauerstoff:}
\begin{align*}
\dot{n} &= D \cdot A \cdot \frac{\Delta c}{l} \\
		&= 0,243 ~\text{cm}^2/\text{s} \cdot (\pi \cdot (1,25 \cdot 10^{-3}~\text{cm})^2) \cdot \frac{0,191 ~\text{cm}^3 \text{/cm$^3$}}{4 \cdot 10^{-3}~\text{cm} } \\
		&= 5,69 \cdot 10^{-5} ~\text{cm$^3$/s}
\end{align*}



\textbf{Wasser:}
\begin{align*}
\dot{n} &= D \cdot A \cdot \frac{\Delta c}{l} \\
		&= 0,243 ~\text{cm}^2/\text{s} \cdot (\pi \cdot (1,25 \cdot 10^{-3}~\text{cm})^2) \cdot \frac{1,291 \cdot 10^{-5} ~\text{g}\text{/cm$^3$}}{4 \cdot 10^{-3}~\text{cm} } \\
		&= 3,85 \cdot 10^{-9} ~\text{g/s}
\end{align*}
}
\pagebreak
\else
\fi



\subsection*{Aufgabe 3}
Ein Lebensmittel ist mit einem sauerstoff- und wasserdampfundurchlässigen Packstoff
verpackt, in dem jedoch eine Pore (Radius: 12,5 µm) ist. Die Pore ist durch eine
Kunststoffschicht mit der Dicke 40 µm abgedeckt. Berechnen Sie den Stoffstrom für
Sauerstoff und Wasserdampf, wenn der Partialdruckunterschied p$_{O2}$ 0,21 bar und der
Partialdruckunterschied p$_{H2O}$ 14,07 mbar beträgt. Der Permeationskoeffizient für Sauerstoff ist mit 628 (cm$^3$ (STP) $\cdot$ 100 µm)/(m$^2$ $\cdot$ d $\cdot$ bar) gegeben. Der Permeationskoeffizient für Wasserdampf beträgt (16,72 g $\cdot$ 100 µm)/(m$^2$ $\cdot$ d $\cdot$ bar).

%\begin{center}
%\begin{tabular}{ll}
%
%Pore Radius: &12,5 µm \\
%Pore Dicke:		&40 µm \\ 
%Partialdruckunterschied Sauerstoff: &0,21 bar \\
%Partialdruckunterschied Wasser: &14,07 mbar \\
%Permeationskoeffizient Sauerstoff: & (628 cm$^3$ $\cdot$ 100 µm)/(m$^2$ $\cdot$ d $\cdot$ bar) \\
%Permeationskoeffizient Wasser: &(16,72 g $\cdot$ 100 µm)/(m$^2$ $\cdot$ d $\cdot$ bar)
%
%\end{tabular}
%
%\end{center}

\ifshowme
\antwort{
\begin{align*}
J &= P \cdot A \cdot \frac{\Delta p}{l}
\end{align*}

\pp 
\textbf{Sauerstoff:}
\begin{align*}
J &= P \cdot A \cdot \frac{\Delta p}{l} \\
 &= \frac{628 \cdot 10^{-6}~\text{m}^3 \cdot 10^{-4} ~\text{m}}{~\text{m}^2 \cdot ~\text{d} \cdot ~\text{bar}} \cdot \pi  \left( 1,25 \cdot 10^{-5}~\text{m} \right)^2  \cdot \frac{0,21 ~\text{bar}}{4 \cdot 10^{-5} ~\text{m}} \\
  &= 1,62 \cdot 10^{-13} ~\text{m}^3 / ~\text{d}
\end{align*}
\pp 
\textbf{Wasser:}
\begin{align*}
J &= P \cdot A \cdot \frac{\Delta p}{l} \\
 &= \frac{16,72 ~\text{g}  \cdot 10^{-4} ~\text{m}}{~\text{m}^2 \cdot ~\text{d} \cdot ~\text{bar}} \cdot \pi  \left( 1,25 \cdot 10^{-5}~\text{m} \right)^2  \cdot \frac{14,07 \cdot 10^{-3} ~\text{bar}}{4 \cdot 10^{-5} ~\text{m}} \\
  &= 2,88 \cdot 10^{-10} ~\text{g} / ~\text{d}
\end{align*}

}
\pagebreak
\else
\fi



\subsection*{Aufgabe 4}
Durch Mischen von Orangensaft (Extraktgehalt: 12 °Brix, Dichte: 1,0465 kg/L, Gesamtsäure: 9
g/L, berechnet als Weinsäure WS), einer Saccharoselösung (c = 65\% w/w) und Wasser soll ein
Orangennektar mit einem Extraktgehalt von 8 °Brix (Dichte: 1,0299 kg/L) und einem
Gesamtsäuregehalt von 5 g/L (berechnet als WS) hergestellt werden. Ein Zusatz von
Säuerungsmitteln ist bei der Herstellung von Fruchtnektaren nicht üblich.

\begin{itemize}


\item Berechnen Sie, wieviel L Orangennektar aus 1 L Saft hergestellt werden können.
\ifshowme
\antwort{
\begin{align*}
\frac{9~\text{g}}{\text{L}} \cdot 1 \text{L} &= \frac{5~\text{g}}{\text{L}} \cdot x  \\
x &= \frac{9~\text{g}}{\text{L}} \cdot \frac{\text{L}}{5~\text{g}} \\
x &= \frac{9}{5} = 1,8~\text{L}
\end{align*}

}
\else
\fi





\pp
\item Berechnen Sie, welche Mengen an Saccharoselösung und Wasser entsprechend zu dem
Saft zugegeben werden müssen, um 1 kg des Getränkes herzustellen. Tipp: Berechnen Sie
dazu zunächst den Massenanteil an Fruchtsaft im Orangennektar.
\ifshowme
\antwort{
\begin{align*}
m_1 &= 1,8 ~\text{L} \cdot 1,0299 ~\text{kg/L} = 1,854 ~\text{kg} \\
m_2 &= 1 ~\text{L} \cdot 1,0465 ~\text{kg/L} = 1,0465 ~\text{kg} \\ 
&\rightarrow \frac{1,0465 ~\text{kg}}{1,854 ~\text{kg}} =  0,5644 ~\text{kg}
\end{align*}
Um 1 Kg Nektar herstellen zu können benötigen wir 0,5644 kg Saft. \par 
Menge Zucker aus Saft berechnen:
\begin{align*}
0,5644 ~\text{kg} \cdot 0,12 = 0,068 ~\text{kg}
\end{align*}
Da wir 1 kg $\cdot $ 8° Brix = 0,08 kg Zucker in unserem Nektar benötigen, fehlen noch:
\begin{align*}
0,08~\text{kg} - 0,068~\text{kg} = 0,012 ~\text{kg}
\end{align*}
Dafür benötigen wir folgende Menge 65\% Saccharose-Lösung:
\begin{align*}
\frac{0,012~\text{kg}}{0,65} = 0,0185 ~\text{kg}
\end{align*}
Und der Rest ist Wasser:
\begin{align*}
(1-0,5644-0,0185)~\text{kg} = 0,4171 ~\text{kg}
\end{align*}
}
\else
\fi



\end{itemize}

\end{document}
